\section{Approximation and limiting cases}

Mathematics often uses exact expressions, but in applications we frequently use approximations. An approximation replaces a complicated expression with a simpler one under specified conditions.

For small \(x\) measured in radians,
\[
\sin x \approx x.
\]
Also, for small \(x\),
\[
(1+x)^n \approx 1+nx.
\]

\begin{remarkbox}
An approximation is useful only together with its condition of validity. Without that condition, the approximation is incomplete information.
\end{remarkbox}

A limiting case is a way of checking an expression by asking what happens when a variable becomes very small, very large, or takes a special value.

For example, consider
\[
f(x)=\frac{x}{1+x}.
\]
If \(x\) is very small, then
\[
f(x)\approx x.
\]
If \(x\) is very large, then
\[
f(x)\approx 1.
\]
These limiting behaviours help us understand the structure of the function.

\begin{summarybox}
Approximations simplify expressions. Limiting cases help us check whether expressions behave sensibly at special scales or special values.
\end{summarybox}
